\begin{tabularx}{\textwidth}{L{29mm}L{48mm}X}
\toprule
\textbf{Elem} & \textbf{Feladat} & \textbf{Megjegyzés} \\
\midrule
\code{try} & A potenciálisan kivételt dobó normál kód helye. & A hibakezelő kódot nem ide, hanem a megfelelő \code{catch} ágba tesszük. \\
\code{catch} & A megadott típusú kivétel elkapása és kezelése. & Több \code{catch} esetén a specifikusabb kivételt kell előre írni. \\
\code{finally} & Mindenképpen lefutó lezáró blokk. & Erőforrás-felszabadításra használható, de Java-ban gyakran kiváltható \code{try-with-resources} szerkezettel. \\
\code{throw} & Kivétel kézi dobása. & Például hibás paraméter esetén saját kivétel dobható. \\
\code{throws} & Metódusfejlécben jelzi, hogy a metódus milyen ellenőrzött kivételt továbbíthat. & A hívónak kezelnie kell, vagy tovább kell jeleznie. \\
Ellenőrzött kivétel & Fordító által kikényszerített kezelés vagy továbbjelzés. & Példa: sok \code{IOException} jellegű hiba. \\
Nem ellenőrzött kivétel & \code{RuntimeException} vagy leszármazottja. & Általában programozási hibát vagy előfeltétel-sértést jelez, például \code{NullPointerException}. \\
\bottomrule
\end{tabularx}
